\chapter{Grafo de dependencias}
\label{sec:grafo-deps}

\section{Catorce crates Rust más dos vecinos}

El repositorio contiene 14 crates Rust, un crate C++ independiente
(\filepath{libgraphmatch/}) y un paquete TypeScript
(\filepath{platform/mcp/}). El cuadro~\ref{tab:crates} los enumera
por capa.

\begin{table}[H]
\centering
\small
\begin{tabular}{@{}lll@{}}
\toprule
Capa & Crate/paquete & Rol \\
\midrule
core       & \filepath{core/graph-engine}        & grafo, DFS, anomaly, interner \\
core       & \filepath{core/matcher-ffi}         & fachada segura a \filepath{libgraphmatch} \\
core       & \filepath{core/gnn}                 & scorer ONNX (v1 estable, v2 gate) \\
\addlinespace
platform   & \filepath{platform/api}             & \code{GraphHunterApi} facade \\
platform   & \filepath{platform/canonical}       & proyección OCSF v1.4 \\
platform   & \filepath{platform/constrained-decode} & guardrail LLM \\
platform   & \filepath{platform/dsl}             & gramática de hipótesis \\
platform   & \filepath{platform/local-llm}       & backends mock/candle \\
platform   & \filepath{platform/parsers}         & \code{LogParser} + impls \\
platform   & \filepath{platform/siem}            & Sentinel / Elastic streaming \\
platform   & \filepath{platform/sources}         & \code{LogSource} + conectores \\
platform   & \filepath{platform/vrl}             & compilador FieldConfig $\to$ VRL \\
\addlinespace
apps       & \filepath{apps/tauri/src-tauri}     & GUI Tauri + commands \\
apps       & \filepath{apps/cli}                 & CLI \\
\addlinespace
vecino C++ & \filepath{libgraphmatch/}           & motor SIMD AVX2 \\
vecino TS  & \filepath{platform/mcp/}            & 54 tools MCP \\
\bottomrule
\end{tabular}
\caption{Crates y paquetes por capa.}
\label{tab:crates}
\end{table}

\section{Dependencias ideales}

La regla enunciada en el capítulo anterior, aplicada al grafo,
produce tres aristas de capa:

\begin{itemize}
    \item \filepath{apps/*} $\to$ \filepath{platform/*}
    \item \filepath{platform/*} $\to$ \filepath{core/*}
    \item \filepath{core/matcher-ffi} $\to$ \filepath{libgraphmatch/}
    (a través del \code{build.rs} que compila el C++ y el único
    bloque \code{extern "C"}).
\end{itemize}

El vecino TypeScript \filepath{platform/mcp/} no depende de crates
Rust en tiempo de compilación: cada tool llama al HTTP API que
expone el Tauri (o un binario dedicado). Lo clasificamos como
\emph{plataforma} por rol, aunque técnicamente es su propio proceso.

\section{Dependencias reales}

Dos aristas rompen la regla hacia abajo en el estado actual:

\begin{itemize}
    \item \filepath{core/graph-engine} $\to$
    \filepath{platform/sources} (referencia por conveniencia al
    tipo \code{AadEventBatch} usado por el GNN bridge).
    \item \filepath{core/graph-engine} $\to$ \filepath{platform/dsl}
    (acoplamiento pragmático introducido al extraer el DSL; el
    parser necesita helpers del core y el core necesita tipos del
    DSL para materializar hipótesis).
\end{itemize}

Son dos aristas de \emph{core a platform}, dirección prohibida por
el modelo. No son ciclos (el grafo sigue siendo un DAG), pero
violan el ideal de shearing. El capítulo~\ref{sec:deudas} discute
cómo eliminarlas en el siguiente ciclo de trabajo.

\section{Ausencia de ciclos}

La verificación se hace fuera de banda con:

\begin{lstlisting}[language=, basicstyle=\ttfamily\small]
cargo metadata --format-version 1 | \
  jq -r '.packages[] | "\(.name) \(.dependencies[].name)"' | \
  ...construcción de DAG y detección de ciclos
\end{lstlisting}

El resultado al cierre de la Fase~4 del rediseño: \textbf{cero
ciclos}. Dos edges ``indeseados'' documentados, pero nada
bidireccional. Es una propiedad que se reverifica en cada PR que
toca \filepath{Cargo.toml}.

\section{Diagrama resumen}

\begin{figure}[H]
\centering
\begin{tikzpicture}[
    node distance=8mm and 14mm,
    crate/.style={draw, rounded corners=2pt, minimum width=24mm,
                  minimum height=6mm, font=\small\ttfamily,
                  fill=gray!5},
    app/.style={crate, fill=accent!15},
    plat/.style={crate, fill=accent!8},
    core/.style={crate, fill=gray!10},
    ext/.style={crate, fill=gray!20, font=\small\ttfamily\itshape},
    dep/.style={-Latex, thick, gray!70},
    bad/.style={-Latex, thick, red!70, dashed},
]

\node[app] (tauri)   {apps/tauri};
\node[app, right=of tauri] (cli) {apps/cli};
\node[app, right=of cli]   (mcp) {platform/mcp};

\node[plat, below=of tauri] (api)     {platform/api};
\node[plat, right=of api]   (siem)    {platform/siem};
\node[plat, right=of siem]  (can)     {platform/canonical};

\node[plat, below=of api]   (sources) {platform/sources};
\node[plat, right=of sources](parsers){platform/parsers};
\node[plat, right=of parsers](dsl)    {platform/dsl};

\node[core, below=of sources](engine) {core/graph-engine};
\node[core, right=of engine](ffi)    {core/matcher-ffi};
\node[core, right=of ffi]  (gnn)    {core/gnn};

\node[ext, below=of ffi]    (libgm)  {libgraphmatch};

\draw[dep] (tauri) -- (api);
\draw[dep] (cli)   -- (api);
\draw[dep] (mcp)   -- (api);
\draw[dep] (api)   -- (siem);
\draw[dep] (api)   -- (can);
\draw[dep] (api)   -- (sources);
\draw[dep] (api)   -- (parsers);
\draw[dep] (api)   -- (dsl);
\draw[dep] (api)   -- (engine);
\draw[dep] (sources) -- (engine);
\draw[dep] (parsers) -- (engine);
\draw[dep] (dsl)     -- (engine);
\draw[dep] (engine)  -- (ffi);
\draw[dep] (ffi)     -- (libgm);
\draw[dep] (engine)  -- (gnn);

\draw[bad] (engine.east) to[bend right=30] (dsl.south);
\draw[bad] (engine) to[bend left=20] (sources);

\end{tikzpicture}
\caption{Grafo de dependencias (simplificado). Las flechas rojas
punteadas son las dos edges ``core $\to$ platform'' que rompen la
regla; capítulo~\ref{sec:deudas} las discute.}
\label{fig:dep-graph}
\end{figure}

El grafo es un DAG; la única bidireccionalidad conceptual
(\filepath{core} $\leftrightarrow$ \filepath{platform}) es la
deuda explícita que se quiere eliminar.
